\section{Verification and Validation}
\label{sec:vv}

We verify and validate the model along the lines set out by
Rossetti~\cite{rossetti2016}: first confirming it was built right, then that it
is the right model of the line.

\subsection{Verification (did we build the model right?)}
We checked the implementation in four ways. First, we traced a single batch
through every block to confirm the routing, the coating bypass for CoreTab, and
the seize and release of the \texttt{cleaningTeam}. Two degenerate runs then gave
the expected behaviour: with zero failures there is no repair scrap, and with one
product there are no changeovers. A conservation check confirms that the tablets
entering equal the tablets bottled plus the tablets scrapped, and a mass balance
confirms that bottles times bottle size never exceed \num{300000} per batch.

\subsection{Validation and calibration (did we build the right model?)}
For validation we compare the model against analytic expectations. Drying carries
the largest required machine-hours over the plan (about \SI{300}{h},
Table~\ref{tab:capacity}), so it should be the binding constraint, ahead of the
cleaning team and packaging. The simulated baseline utilisation confirms this
ordering: drying \SI{94.9}{\percent}, cleaning team \SI{79.6}{\percent} and
packaging \SI{52.2}{\percent}. The per-station mean processing times also
reproduce the input distributions they were drawn from. The one free modelling
choice, namely how leftover tablets that cannot complete a bottle are rounded
down, was fixed so that the baseline yield and throughput stay internally
consistent before we ran the experiments.

\begin{table}[H]
\centering
\small
\caption{Verification and validation checks against the $n=20$ baseline.}
\label{tab:verify}
\begin{tabularx}{\textwidth}{X l l c}
\toprule
Check & Expected & Simulated (baseline) & OK \\
\midrule
Batch conservation   & in $=$ shipped $+$ scrapped $=50$ & $46.9 + 3.1 = 50.0$ & \checkmark \\
Primary constraint   & drying (largest unit-h)           & drying \SI{94.9}{\percent} util & \checkmark \\
Constraint ordering  & drying $>$ clean $>$ packaging     & $94.9 > 79.6 > 52.2$ & \checkmark \\
Little's Law         & $\overline{\text{WIP}} \approx \lambda\overline{W}$ & $19.36$ vs.\ $19.53$ (\SI{0.9}{\percent}) & \checkmark \\
\bottomrule
\end{tabularx}
\end{table}

\subsection{Analytical capacity check (predict, then confirm)}
\label{sec:capacity}
Before trusting the simulation we predict the bottleneck from first principles:
the required machine-hours each operation must perform over the 50-batch plan.

\begin{table}[H]
\centering
\caption{Required machine-hours over the plan (predicts the constraint).}
\label{tab:capacity}
\begin{tabularx}{\textwidth}{l c c c r}
\toprule
Operation & h/batch (incl.\ clean) & Units & Batches & Required unit-h \\
\midrule
Premix (old$\parallel$new)   & $\approx$4.75 & 2 & 50 & $\approx$119 \\
Granulation                  & 1.0  & 1 & 50 & 50 \\
\textbf{Drying} (+2 h clean) & $\approx$6.0 & 1 & 50 & \textbf{$\approx$300} \\
Compression                  & 2.0  & 3 & 50 & 33 \\
Coating (+1 h clean)         & 3.0  & 1 & 33 & 99 \\
Packaging                    & $\approx$9--10 & 3 & 50 & $\approx$150--167 \\
\textbf{Cleaning team} (shared) & - & 1 & - & \textbf{$\approx$170} \\
\bottomrule
\end{tabularx}
\end{table}

The analysis predicts \textbf{drying} as the primary constraint, the
\textbf{cleaning team} second, and packaging third. The 2-hour drying clean
couples the top two, since it both blocks the dryer and is the largest single
draw on the cleaning team. Section~\ref{sec:expresults} confirms this against
the simulated utilisation.

\subsection{Little's Law consistency check}
\label{sec:little}
As a further validation we verify Little's Law on the baseline,
$\text{WIP} = \lambda \cdot W$, linking our three flow KPIs:
\begin{equation}
  \overline{\text{WIP}} \;\approx\; \text{throughput rate} \times \overline{W},
  \label{eq:little}
\end{equation}
where $\overline{W}$ is mean flow time. The throughput rate is taken
\emph{independently} of the WIP and flow-time statistics, from the shipped-batch
count over the makespan: $\lambda = 46.9/436.2 = 0.1075$ batches/h. With
$\overline{W}=\SI{180.0}{\hour}$ this
predicts $\lambda\overline{W}=19.36$ batches, against a directly measured
$\overline{\text{WIP}}=19.53$ batches, a \SI{0.9}{\percent} discrepancy that is
well within sampling error. The three flow KPIs are therefore mutually
consistent, which confirms that the model conserves work.
